\chapter{Linear Models under Stress}\label{ch:qm:linear-models-under-stress}

A factor model regresses a stock's daily returns on five factors, month by month. Two of the factors move together, with
correlation 0.97. Between the first month and the second, the stock's beta to the first factor goes from 0.52 to $-1.59$ and its
beta to the second from 0.84 to 2.76; their sum moves only from 1.36 to 1.17, and the second month's fitted returns, recomputed
with the first month's betas, keep a correlation of 0.92 with its own fit. The regression is not wrong: it has been asked a
question, how the exposure splits between two nearly identical factors, that three weeks of data cannot answer. This chapter is
about least squares when its assumptions strain: the geometry that explains it, collinearity, the penalties that stabilise it,
regressors measured with error, and the \omterm{def:qm:estimation:estimator}{standard errors} of regressions run on panels of stocks and months.

\section{The geometry of least squares}

\begin{definition}[Ordinary least squares, hat matrix]\label{def:qm:linear-models-under-stress:ols}
For $y \in \R^n$ and a full-rank design $X \in \R^{n \times k}$, \emph{ordinary least squares}\index{ordinary least squares} (OLS) chooses
$\hat\beta = (X^\top X)^{-1}X^\top y$, the minimiser of $\lvert y - X\beta\rvert^2$. The fitted values are $\hat y = Hy$ with the \emph{hat
matrix}\index{hat matrix} $H = X(X^\top X)^{-1}X^\top$, the orthogonal projection onto the column space of $X$; its diagonal entries
$h_{ii}$, which sum to $k$, are the leverages of the observations.
\end{definition}

The residual $y - \hat y = (I - H)y$ is orthogonal to every column of $X$. An observation with leverage near one pulls the fit
through itself: one bad print in a regressor is the worst kind of outlier, because it is far from the others in the direction the
regression looks.

\begin{theorem}[Gauss--Markov]\label{thm:qm:linear-models-under-stress:gm}
If $y = X\beta + \varepsilon$ with $\E[\varepsilon \mid X] = 0$ and $\Var(\varepsilon \mid X) = \sigma^2I$, then $\hat\beta$ is unbiased with
variance $\sigma^2(X^\top X)^{-1}$, and every other linear \omterm{def:qm:estimation:estimator}{unbiased estimator} $Cy$ has variance larger by a positive semidefinite matrix.
\end{theorem}

\begin{proof}
Write $C = (X^\top X)^{-1}X^\top + D$. Unbiasedness for every $\beta$ forces $DX = 0$. Then $\Var(Cy) = \sigma^2CC^\top = \sigma^2(X^\top X)^{-1} +
\sigma^2DD^\top$, the cross terms vanishing because $DX = 0$.
\end{proof}

The theorem is a statement about \omterm{def:qm:estimation:estimator}{unbiased estimators}; the rest of the chapter is about why a biased one is often better, and about
what happens when $\Var(\varepsilon) \ne \sigma^2I$ or the regressor itself is noisy.

\begin{theorem}[Frisch--Waugh--Lovell]\label{thm:qm:linear-models-under-stress:fwl}
In the regression of $y$ on $(X_1, X_2)$, the coefficient on $X_2$ equals the coefficient from regressing $M_1y$ on $M_1X_2$, where $M_1 =
I - X_1(X_1^\top X_1)^{-1}X_1^\top$ removes the part explained by $X_1$; the residuals of the two regressions coincide.
\end{theorem}

\begin{proof}
From $y = X_1\hat\beta_1 + X_2\hat\beta_2 + e$ with $e$ orthogonal to both blocks, apply $M_1$: $M_1y = M_1X_2\hat\beta_2 + e$, since $M_1X_1 = 0$ and
$M_1e = e$. As $e$ is orthogonal to $M_1X_2$, this is the least-squares decomposition of $M_1y$ on $M_1X_2$.
\end{proof}

A beta ``controlling for'' the market is the beta of the stock's market-residual return on the factor's market-residual return. The
theorem is also the fastest way to see collinearity: the coefficient on $X_2$ is estimated from what is left of $X_2$ after $X_1$, and
when little is left, the estimate rests on little.

\section{Collinearity}

\begin{definition}[Multicollinearity, variance inflation factor]\label{def:qm:linear-models-under-stress:vif}
\emph{Multicollinearity}\index{multicollinearity} is near-linear dependence among the columns of the design. The \emph{variance inflation
factor}\index{variance inflation factor} of regressor $j$ is $\mathrm{VIF}_j = 1/(1 - R_j^2)$, with $R_j^2$ the $R^2$ of regressing $x_j$ on the
other regressors.
\end{definition}

By the Frisch--Waugh--Lovell theorem, $\Var(\hat\beta_j) = \sigma^2/(n\,\widehat{\Var}(x_j)(1 - R_j^2))$: the variance of a coefficient is the
variance it would have alone, multiplied by its VIF. Two factors with correlation 0.97 have $\mathrm{VIF} = 1/(1 - 0.97^2) = 16.9$ (18.5 in
the chapter's sample), so each beta is four times noisier than it would be alone; the sum of the two betas, the exposure to their common
direction, is not inflated at all.

Over 24 simulated months of 21 days, the stock's true betas to the two collinear factors are 0.6 and 0.4. The monthly estimates have standard
deviations of 1.21 and 1.23 across months, a sign opposite to the truth in 6 and 8 months, while their sum has a standard deviation of 0.23
(\cref{fig:qm:linear-models-under-stress:betas}); the betas to the three uncorrelated factors vary by 0.26 to 0.32. A risk report that lists
the two betas separately reports noise; one that reports their sum, or projects the two factors on their first principal component,
reports the exposure.

\begin{omfigure}
\begin{tikzpicture}
\begin{axis}[width=11cm,height=5cm,xlabel={month},ylabel={monthly OLS beta},
  tick label style={font=\small},label style={font=\small},xmin=0.5,xmax=24.5,ymin=-3.5,ymax=4,grid=major,grid style={gray!25},
  legend columns=3,legend style={font=\footnotesize,at={(0.5,-0.3)},anchor=north,draw=none,fill=none,
  /tikz/every even column/.append style={column sep=8pt}},table/col sep=comma]
\addplot[omDef,thick,mark=*,mark size=1.3pt] table[x=month,y=b1]{figdata/methods/16-linear-models-under-stress/betas.csv};
\addplot[omThm,thick,mark=square*,mark size=1.3pt] table[x=month,y=b2]{figdata/methods/16-linear-models-under-stress/betas.csv};
\addplot[black,very thick] table[x=month,y=sum]{figdata/methods/16-linear-models-under-stress/betas.csv};
\legend{beta to factor 1,beta to factor 2,sum of the two}
\end{axis}
\end{tikzpicture}

\omcaption{A stock's betas to two factors with correlation 0.97, estimated each month from 21 daily returns (true values 0.6 and 0.4). The
individual betas swing between $-2.7$ and $3.3$; their sum stays near its true value of 1.0 (standard deviation 0.23). Data: the chapter's
tutorial, seeded.}\label{fig:qm:linear-models-under-stress:betas}
\end{omfigure}

\section{Ridge, lasso and elastic net}

\begin{definition}[Regularisation, ridge, lasso, elastic net]\label{def:qm:linear-models-under-stress:penalties}
\emph{Regularisation}\index{regularisation} adds a penalty on the size of the coefficients to the least-squares objective. \emph{Ridge
regression}\index{ridge regression} minimises $\lvert y - X\beta\rvert^2 + \lambda\lvert\beta\rvert_2^2$; the \emph{lasso}\index{lasso} minimises
$\frac1{2n}\lvert y - X\beta\rvert^2 + \lambda\lvert\beta\rvert_1$; the \emph{elastic net}\index{elastic net} uses the penalty $\lambda\bigl(\alpha\lvert\beta\rvert_1
+ \frac{1 - \alpha}2\lvert\beta\rvert_2^2\bigr)$, $\alpha \in [0, 1]$.
\end{definition}

\begin{proposition}[Ridge shrinks along the singular directions]\label{prop:qm:linear-models-under-stress:ridge}
If $X = UDV^\top$ is a thin singular value decomposition, the ridge estimate is $\hat\beta_\lambda = \sum_jv_j\frac{d_j}{d_j^2 + \lambda}u_j^\top y$:
the least-squares coefficient along each right singular vector $v_j$ is multiplied by $d_j^2/(d_j^2 + \lambda)$.
\end{proposition}

\begin{proof}
$\hat\beta_\lambda = (X^\top X + \lambda I)^{-1}X^\top y = V(D^2 + \lambda I)^{-1}DU^\top y$, and least squares is the case $\lambda = 0$, with factor
$1/d_j$ along $v_j$.
\end{proof}

Collinear directions have small singular values, and they are the ones ridge shrinks: for the two factors above, the difference $F_1 - F_2$ has
a small singular value and its coefficient is pulled toward zero, while the sum is left almost alone. Ridge is the posterior mean under a normal
prior on $\beta$ (chapter~14). The \omterm{def:qm:linear-models-under-stress:penalties}{lasso}'s absolute-value penalty has a corner at zero: for an orthonormal design its solution is the
least-squares coefficient soft-thresholded, $\operatorname{sign}(b)\max(|b| - \lambda, 0)$, so it sets small coefficients exactly to zero and
selects variables. Within a block of correlated predictors it tends to pick one and drop the others, somewhat arbitrarily; the \omterm{def:qm:linear-models-under-stress:penalties}{elastic net}'s
ridge part spreads the weight across the block. The \omterm{def:qm:linear-models-under-stress:penalties}{lasso} and \omterm{def:qm:linear-models-under-stress:penalties}{elastic net} are computed by cyclic coordinate descent, each coordinate being a
one-dimensional soft-thresholding (Friedman, Hastie and Tibshirani, 2010).

\begin{definition}[Cross-validation]\label{def:qm:linear-models-under-stress:cv}
\emph{Cross-validation}\index{cross-validation} chooses a penalty by fitting on part of the data and measuring the prediction error on the rest,
over several splits. With time series, every training set must precede its test set, with a gap (a purge) at least as long as the dependence
between the targets, so that no information from the test period leaks into the fit.
\end{definition}

The tutorial predicts a daily return from thirty predictors in three blocks of ten, correlated 0.8 within each block; five predictors carry
signal, and the signal is weak (a perfect model would explain 5.0\% of the variance out of sample). The first 1\,500 days are split into
five expanding folds with a purge of five days, and the last 500 are the test. OLS on all thirty predictors explains 1.5\% of the test
variance; ridge at the cross-validated penalty 3.5\%, the \omterm{def:qm:linear-models-under-stress:penalties}{lasso} 3.4\%, the \omterm{def:qm:linear-models-under-stress:penalties}{elastic net} 3.2\%. The \omterm{def:qm:linear-models-under-stress:penalties}{lasso} keeps five predictors, three of them
right; in the two blocks where it errs it has kept a correlated neighbour of a true predictor
(\cref{fig:qm:linear-models-under-stress:path,fig:qm:linear-models-under-stress:cv}).

\begin{omfigure}
\begin{tikzpicture}
\begin{semilogxaxis}[width=11cm,height=5.4cm,xlabel={lasso penalty $\lambda$},ylabel={coefficient (standardised)},x dir=reverse,
  tick label style={font=\small},label style={font=\small},xmin=1e-4,xmax=1,ymin=-0.2,ymax=0.2,grid=major,grid style={gray!25},
  table/col sep=comma]
\pgfplotsinvokeforeach{1,2,4,5,6,7,8,9,10,13,14,15,16,17,18,19,20,21,22,23,24,26,27,28,29}{
  \addplot[black!35,thin] table[x=lam,y=b#1]{figdata/methods/16-linear-models-under-stress/path.csv};
}
\pgfplotsinvokeforeach{0,3,11,12,25}{
  \addplot[omDef,very thick] table[x=lam,y=b#1]{figdata/methods/16-linear-models-under-stress/path.csv};
}
\addplot[omThm,dashed,thick,sharp plot] coordinates {(0.0398,-0.2) (0.0398,0.2)};
\end{semilogxaxis}
\end{tikzpicture}

\omcaption{\omterm{def:qm:linear-models-under-stress:penalties}{Lasso} path of the thirty standardised predictors as the penalty falls from 1 to $10^{-4}$: the five true signals in blue, the
twenty-five noise predictors in grey, and the cross-validated penalty 0.040 (dashed), where five coefficients are nonzero. Data: the
chapter's tutorial, seeded.}\label{fig:qm:linear-models-under-stress:path}
\end{omfigure}

\begin{omfigure}
\begin{tikzpicture}
\begin{semilogxaxis}[width=10cm,height=5cm,xlabel={penalty $\lambda$},ylabel={cross-validated MSE},x dir=reverse,
  tick label style={font=\small},label style={font=\small},xmin=1e-4,xmax=1,grid=major,grid style={gray!25},
  yticklabel style={/pgf/number format/fixed,/pgf/number format/precision=3},
  legend columns=3,legend style={font=\footnotesize,at={(0.5,-0.3)},anchor=north,draw=none,fill=none,
  /tikz/every even column/.append style={column sep=8pt}},table/col sep=comma]
\addplot[omDef,very thick] table[x=lam,y=ridge]{figdata/methods/16-linear-models-under-stress/cv.csv};
\addplot[omThm,very thick,dashed] table[x=lam,y=lasso]{figdata/methods/16-linear-models-under-stress/cv.csv};
\addplot[black,thick,dotted] table[x=lam,y=enet]{figdata/methods/16-linear-models-under-stress/cv.csv};
\legend{ridge,lasso,{elastic net, $\alpha = 0.5$}}
\end{semilogxaxis}
\end{tikzpicture}

\omcaption{Time-ordered \omterm{def:qm:linear-models-under-stress:cv}{cross-validation} error of the three penalised regressions against the penalty (five expanding folds with a five-day
purge, on the first 1\,500 days). Each curve has an interior minimum: too little penalty fits noise, too much discards signal. Data: the
chapter's tutorial, seeded.}\label{fig:qm:linear-models-under-stress:cv}
\end{omfigure}

\section{Weighted and errors-in-variables regression}

\begin{definition}[Weighted and generalised least squares]\label{def:qm:linear-models-under-stress:gls}
\emph{Weighted least squares}\index{weighted least squares} minimises $\sum_iw_i(y_i - x_i^\top\beta)^2$, with weights inversely proportional to the
variance of each observation's error. \emph{Generalised least squares}\index{generalised least squares} minimises $(y - X\beta)^\top\Omega^{-1}(y -
X\beta)$ for an error covariance $\Omega$; it is OLS after multiplying the model by $\Omega^{-1/2}$, and it is the efficient linear \omterm{def:qm:estimation:estimator}{unbiased
estimator} when $\Omega$ is known.
\end{definition}

When $\Omega$ is unknown, the usual choice in finance is OLS with a \omterm{def:qm:estimation:estimator}{standard error} that does not assume it away: the sandwich of
chapter~11, HAC for time series, clustered for panels (below).

\begin{definition}[Errors-in-variables, attenuation bias, total least squares]\label{def:qm:linear-models-under-stress:eiv}
An \emph{errors-in-variables}\index{errors-in-variables} model observes the regressor with noise, $\tilde x = x + u$, with $u$ independent of $x$
and of the equation error. The resulting shrinkage of the OLS slope toward zero is the \emph{attenuation bias}\index{attenuation bias}. \emph{Total
least squares}\index{total least squares} fits the line minimising the sum of squared perpendicular distances, which is consistent when the errors in
$x$ and in $y$ have equal variances.
\end{definition}

\begin{proposition}[The attenuation factor]\label{prop:qm:linear-models-under-stress:attenuation}
If $y = \beta x + \varepsilon$ and $\tilde x = x + u$ with $x$, $u$, $\varepsilon$ independent, the OLS slope of $y$ on $\tilde x$ converges to
$\lambda\beta$ with $\lambda = \Var(x)/(\Var(x) + \Var(u))$.
\end{proposition}

\begin{proof}
$\Cov(y, \tilde x) = \beta\Var(x)$ and $\Var(\tilde x) = \Var(x) + \Var(u)$; the slope is their ratio.
\end{proof}

A desk hedges a bond with a bond future, and estimates the hedge ratio by regressing the bond's daily price changes on the future's. The bond
is marked at a synchronous mid at the close; the future's recorded daily price is its last trade before the close, stale by a few minutes and
at bid or ask. Model the recorded level as the true level plus an independent error of standard deviation 0.14 points; the future's true
daily change has standard deviation 0.40, and the bond moves 0.85 times the future plus an idiosyncratic 0.05. The recorded change carries
the difference of two level errors, of variance $2 \times 0.14^2 = 0.039$, so $\lambda = 0.16/(0.16 + 0.039) = 0.80$. On two years of seeded data
the OLS hedge ratio is 0.68, and across 2\,000 simulated histories it averages 0.683 with a standard deviation of 0.017: precisely wrong.

Two corrections come from the structure of the error. A level error makes consecutive recorded changes negatively correlated, with first
autocovariance $-\sigma_e^2$ (Roll's 1984 argument for the bid-ask spread), so $\hat\lambda = (\hat\gamma(0) + 2\hat\gamma(1))/\hat\gamma(0)$
estimates the attenuation from the future's own series; here the first autocorrelation is $-0.098$, $\hat\lambda = 0.80$, and the
corrected ratio is 0.85. Across histories this correction is unbiased (0.862 on average) but noisy (standard deviation 0.095). Alternatively,
weekly changes add five days of true movement to the same two level errors, so their attenuation is only $5 \times 0.16/(5 \times 0.16 + 0.039) =
0.95$: the weekly regression gives 0.80 on this history, 0.810 on average with a standard deviation of 0.022
(\cref{fig:qm:linear-models-under-stress:hedge}). \omterm{def:qm:linear-models-under-stress:eiv}{Total least squares} gives 0.74: its assumption of equal error variances is false here.

\begin{omfigure}
\begin{tikzpicture}
\begin{axis}[width=10cm,height=5cm,xlabel={estimated hedge ratio},ylabel={density over 2\,000 histories},
  tick label style={font=\small},label style={font=\small},xmin=0.5,xmax=1.2,ymin=0,ymax=26,grid=major,grid style={gray!25},
  legend columns=4,legend style={font=\footnotesize,at={(0.5,-0.3)},anchor=north,draw=none,fill=none,
  /tikz/every even column/.append style={column sep=6pt}},table/col sep=comma]
\addplot[omThm,thick,const plot mark mid,fill=omThm!15] table[x=h,y=ols]{figdata/methods/16-linear-models-under-stress/hedge.csv};
\addplot[black,thick,const plot mark mid] table[x=h,y=roll]{figdata/methods/16-linear-models-under-stress/hedge.csv};
\addplot[omDef,very thick,const plot mark mid] table[x=h,y=weekly]{figdata/methods/16-linear-models-under-stress/hedge.csv};
\addplot[black,dashed,thin,sharp plot] coordinates {(0.85,0) (0.85,26)};
\legend{daily OLS,Roll-corrected,weekly OLS,true 0.85}
\end{axis}
\end{tikzpicture}

\omcaption{Sampling distributions of three hedge-ratio estimates over 2\,000 simulated two-year histories: daily OLS on the noisy future prices
(mean 0.683, attenuated by 0.80), the same divided by the attenuation estimated from the future's first autocorrelation (mean 0.862, standard
deviation 0.095), and OLS on weekly changes (mean 0.810, standard deviation 0.022). Data: the chapter's tutorial, seeded.}\label{fig:qm:linear-models-under-stress:hedge}
\end{omfigure}

The cost is in the residual. Measured against the true synchronous changes, the hedge at the true ratio leaves a daily residual of 0.049
points; at the daily OLS ratio 0.086, 74\% more; at the corrected ratio 0.049; at the weekly ratio 0.053. Unhedged, the bond moves 0.34 a day.

\section{Cross-sectional regressions over time, and panels}

\begin{definition}[Panel data, fixed effects]\label{def:qm:linear-models-under-stress:panel}
\emph{Panel data}\index{panel data} observe many units (stocks, funds, counterparties) over many periods. A \emph{fixed effects}\index{fixed effects}
regression gives each unit (or period) its own intercept, which by the Frisch--Waugh--Lovell theorem is the same as demeaning every variable
within the unit.
\end{definition}

\begin{definition}[Fama--MacBeth regression, clustered standard errors]\label{def:qm:linear-models-under-stress:fm}
A \emph{Fama--MacBeth regression}\index{Fama--MacBeth regression} (Fama and MacBeth, 1973) runs one cross-sectional regression per period and
reports the time-series mean of the slopes, with \omterm{def:qm:estimation:estimator}{standard error} their standard deviation over $\sqrt T$. \emph{Clustered standard
errors}\index{clustered standard errors} compute the \omterm{def:qm:estimation:sandwich}{sandwich variance} with the scores summed within clusters (firms, or periods), allowing any
correlation inside a cluster and none across clusters.
\end{definition}

The Fama--MacBeth \omterm{def:qm:estimation:estimator}{standard error} is right when the slopes are independent over time; a period effect (all stocks' residuals moving together in
a month) is handled, since each month's slope absorbs it. A persistent firm effect is not: if both a characteristic and the residual have a
component fixed for each firm, every month's slope shares the same error, which the time-series standard deviation cannot see and a Newey--West
correction with a few lags cannot see either (Petersen, 2009). In a simulated panel of 100 firms over 60 months with firm effects making up half
the variance of the characteristic and of the residual, the slope's true standard deviation across 500 panels is 0.105. Classical and White
\omterm{def:qm:estimation:estimator}{standard errors} average 0.026, the Fama--MacBeth one 0.023 (0.022 with three Newey--West lags), and only the \omterm{def:qm:estimation:estimator}{standard error} clustered by firm, 0.101,
is right (\cref{fig:qm:linear-models-under-stress:panel}).

\begin{omfigure}
\begin{tikzpicture}
\begin{axis}[width=10cm,height=4.8cm,xbar,bar width=8pt,
  xlabel={standard error of the slope (average over 500 panels)},
  ytick={0,1,2,3,4,5},yticklabels={true standard deviation,classical OLS,White,clustered by firm,Fama--MacBeth,{Fama--MacBeth, 3 lags}},
  tick label style={font=\small},label style={font=\small},
  xmin=0,xmax=0.13,ymin=-0.6,ymax=5.6,grid=major,grid style={gray!25},
  xticklabel style={/pgf/number format/fixed},scaled x ticks=false,table/col sep=comma]
\addplot[fill=omDef!60,draw=omDef,nodes near coords,nodes near coords style={font=\footnotesize,/pgf/number format/fixed,/pgf/number format/precision=3}]
  table[x=se,y=k]{figdata/methods/16-linear-models-under-stress/panel.csv};
\end{axis}
\end{tikzpicture}

\omcaption{\omterm{def:qm:estimation:estimator}{Standard errors} of a panel slope when the characteristic and the residual both have a persistent firm component: only clustering by
firm recovers the true sampling standard deviation (0.105); the classical, White and Fama--MacBeth \omterm{def:qm:estimation:estimator}{standard errors} understate it about fourfold.
Data: the chapter's tutorial, seeded.}\label{fig:qm:linear-models-under-stress:panel}
\end{omfigure}

\section{Tutorial: the hedge ratio that shrank}\label{tut:qm:linear-models-under-stress:hedge}

\begin{tutorial}
\textbf{Goal.} Stabilise a collinear regression, select predictors with time-ordered \omterm{def:qm:linear-models-under-stress:cv}{cross-validation}, and repair an attenuated hedge ratio.
\textbf{End state:} \cref{fig:qm:linear-models-under-stress:path,fig:qm:linear-models-under-stress:cv,fig:qm:linear-models-under-stress:hedge}, the
test $R^2$ of 1.5\% (OLS) and 3.5\% (ridge), and the hedge ratios 0.68, 0.85 and 0.80.
\begin{tutsteps}
\item \textbf{Ridge through the SVD and the \omterm{def:qm:linear-models-under-stress:penalties}{elastic net} by coordinate descent.}
\omcode{code/firm/linreg/firm_linreg.py}{76}{104}{Ridge and elastic net.}\label{lst:qm:linear-models-under-stress:enet}
\item \textbf{Time-ordered folds with a purge.}
\omcode{code/firm/linreg/firm_linreg.py}{107}{117}{Expanding-window cross-validation folds.}\label{lst:qm:linear-models-under-stress:folds}
\item \textbf{Run} \texttt{fit\_all()}, \texttt{hedge()}, \texttt{hedge\_mc()} and \texttt{panel\_truth()} in \texttt{qm\_linreg.py}, then
\texttt{fig\_linreg.py}.
\end{tutsteps}
\textbf{What to change next.} Replace the two collinear factors by their sum and difference and watch the difference's beta alone carry the
noise; shorten the purge to zero with overlapping five-day targets and measure how much \omterm{def:qm:linear-models-under-stress:cv}{cross-validation} flatters the fit.
\end{tutorial}

\section{Build: the regression module}\label{bld:qm:linear-models-under-stress:linreg}

\begin{build}
\bfield{Purpose} Every regression the miniature firm runs (betas, hedge ratios, signal fits, panel tests) goes through one module whose
\omterm{def:qm:estimation:estimator}{standard errors} are chosen, not defaulted.
\bfield{Interface} \texttt{ols(X, y, cov, lags, groups)} with \texttt{cov} in classical, \texttt{hc}, \texttt{hac}, \texttt{cluster};
\texttt{vif(X)}; \texttt{ridge(X, y, lam)}; \texttt{elastic\_net(X, y, lam, alpha)}; \texttt{time\_series\_folds(n, n\_folds, purge)};
\texttt{fama\_macbeth(ys, Xs, lags)}; \texttt{tls(x, y)}.
\bfield{Rules} The covariance type is a required choice in reports; predictors are standardised on the training fold only; folds never train on
data after their test block; warm starts along a penalty path.
\bfield{Acceptance tests} \texttt{code/firm/linreg/tests/}: White \omterm{def:qm:estimation:estimator}{standard errors} match the simulated sampling spread under heteroskedasticity;
the Frisch--Waugh--Lovell identity; the VIF of two variables with correlation 0.97; ridge equals its normal equations and the \omterm{def:qm:linear-models-under-stress:penalties}{elastic net} with
$\alpha = 0$ equals ridge; the \omterm{def:qm:linear-models-under-stress:penalties}{lasso} soft-thresholds an orthonormal design; purged folds; Fama--MacBeth on independent periods.
\bfield{Stretch} Deming regression with a known error-variance ratio; the adaptive \omterm{def:qm:linear-models-under-stress:penalties}{lasso}; two-way clustering by firm and month.
\end{build}

\begin{omsources}
\item R.~Frisch and F.~V.~Waugh, ``Partial time regressions as compared with individual trends'', \emph{Econometrica} 1, 1933; M.~C.~Lovell,
\emph{Journal of the American Statistical Association} 58, 1963.
\item A.~E.~Hoerl and R.~W.~Kennard, ``Ridge regression: biased estimation for nonorthogonal problems'', \emph{Technometrics} 12, 1970.
\item R.~Tibshirani, ``Regression shrinkage and selection via the lasso'', \emph{Journal of the Royal Statistical Society B} 58, 1996.
\item H.~Zou and T.~Hastie, ``Regularization and variable selection via the elastic net'', \emph{Journal of the Royal Statistical Society B}
67, 2005.
\item J.~Friedman, T.~Hastie and R.~Tibshirani, ``Regularization paths for generalized linear models via coordinate descent'', \emph{Journal of
Statistical Software} 33, 2010.
\item E.~F.~Fama and J.~D.~MacBeth, ``Risk, return, and equilibrium: empirical tests'', \emph{Journal of Political Economy} 81, 1973.
\item M.~A.~Petersen, ``Estimating standard errors in finance panel data sets: comparing approaches'', \emph{Review of Financial Studies} 22, 2009.
\item R.~Roll, ``A simple implicit measure of the effective bid-ask spread in an efficient market'', \emph{Journal of Finance} 39, 1984.
\end{omsources}

\section{Exercises}

\begin{exercise}[$\star$]\label{exo:qm:linear-models-under-stress:1}
Two regressors have correlation 0.9. What is the \omterm{def:qm:linear-models-under-stress:vif}{variance inflation factor}, and by what factor does each coefficient's \omterm{def:qm:estimation:estimator}{standard error} grow?
\end{exercise}

\begin{exercise}[$\star$]\label{exo:qm:linear-models-under-stress:2}
A regressor's true variance is 1 and it is measured with independent noise of variance 0.25. If the true slope is 2, what does OLS estimate?
\end{exercise}

\begin{exercise}[$\star$]\label{exo:qm:linear-models-under-stress:3}
Show that the leverages $h_{ii}$ sum to $k$, the number of columns of $X$.
\end{exercise}

\begin{exercise}[$\star\star$]\label{exo:qm:linear-models-under-stress:4}
For an orthonormal design ($X^\top X = nI$), derive the \omterm{def:qm:linear-models-under-stress:penalties}{lasso} solution coordinate by coordinate.
\end{exercise}

\begin{exercise}[$\star\star$]\label{exo:qm:linear-models-under-stress:5}
A future's recorded daily changes have variance 0.199 and first autocovariance $-0.0196$. Estimate the level-error variance and the attenuation
factor of a daily hedge regression, and the attenuation of a weekly one.
\end{exercise}

\begin{exercise}[$\star\star$]\label{exo:qm:linear-models-under-stress:6}
With the two collinear factors of the chapter ($\rho = 0.97$, unit variances), which linear combination of the two betas does the data pin down,
and how much more precisely than each beta?
\end{exercise}

\begin{exercise}[$\star\star\star$]\label{exo:qm:linear-models-under-stress:7}
\emph{Coding.} Simulate the chapter's panel with firm effects and compare the average classical, clustered and Fama--MacBeth \omterm{def:qm:estimation:estimator}{standard errors} with
the true sampling standard deviation of the slope over 500 panels.
\end{exercise}

\begin{exercise}[$\star\star\star$]\label{exo:qm:linear-models-under-stress:8}
\emph{Find the flaw.} ``We chose the \omterm{def:qm:linear-models-under-stress:penalties}{lasso} penalty by ten-fold \omterm{def:qm:linear-models-under-stress:cv}{cross-validation} with random folds on five years of daily data, predicting each
stock's next twenty-day return; the cross-validated $R^2$ is 8\%.''
\end{exercise}

\section{Problem: The Hedge Ratio That Shrank}

\begin{problem}[{Weekend problem --- hedging a bond with a future whose prices are noisy}]\label{pb:qm:linear-models-under-stress:1}
A bond's daily price change is 0.85 times the future's true change plus an idiosyncratic 0.05 points; the future's true change has standard
deviation 0.40 points; its recorded daily price carries an independent level error of standard deviation 0.14. Two years (500 days) of the
chapter's seeded data are available.

\medskip\noindent\textbf{Part I --- The naive hedge.}
\begin{enumerate}
\item What hedge ratio does daily OLS give?
\item What attenuation factor does the model predict, and what does it imply for the average OLS ratio?
\item Across 2\,000 simulated histories, what are the mean and standard deviation of the daily OLS ratio?
\item Why is a precise estimate not a correct one here?
\item What residual daily risk does the OLS ratio leave, against the true ratio's?
\end{enumerate}

\medskip\noindent\textbf{Part II --- Corrections.}
\begin{enumerate}[resume]
\item What is the first autocorrelation of the recorded future changes, and what does it imply for $\hat\lambda$?
\item What is the corrected ratio, and its residual risk?
\item What attenuation does a weekly regression suffer, and what ratio does it give?
\item What are the mean and standard deviation of the corrected and weekly ratios across histories?
\item What does \omterm{def:qm:linear-models-under-stress:eiv}{total least squares} give, and why is it wrong here?
\end{enumerate}

\medskip\noindent\textbf{Part III --- Beyond the hedge.}
\begin{enumerate}[resume]
\item What would the attenuation be with a level error of 0.07 instead of 0.14?
\item How would synchronous prices (both marked at the same instant, at mid) change the problem?
\item If the hedge ratio were estimated on monthly changes, what would be gained and lost?
\item Could the collinearity of the chapter's factor model be repaired by the same means?
\item How would you check the estimated ratio before trading it?
\end{enumerate}

\medskip\noindent\textbf{Part IV --- Judgement.}
\begin{enumerate}[resume]
\item Which ratio would you use, and why?
\item What should the data team change?
\item When is the attenuation of a regression slope harmless?
\item State the \emph{named result}: the attenuation factor, the corrected ratio, and the residual risk each leaves.
\item In one sentence: what does noise in a regressor do that noise in the dependent variable does not?
\end{enumerate}
\end{problem}

\section{Interview questions}

\begin{interviewq}[$\star$ \iqroles{researcher, mle}]\label{iq:qm:linear-models-under-stress:1}
Two of your regressors have correlation 0.97. What happens to their coefficients, and what do you do?
\end{interviewq}

\begin{interviewq}[$\star\star$ \iqroles{researcher, trader}]\label{iq:qm:linear-models-under-stress:2}
Your regression hedge ratio is 0.68 but the duration ratio says 0.85. What could explain the gap?
\end{interviewq}

\begin{interviewq}[$\star\star$ \iqroles{mle, researcher}]\label{iq:qm:linear-models-under-stress:3}
Ridge or \omterm{def:qm:linear-models-under-stress:penalties}{lasso}: what does each do to correlated predictors?
\end{interviewq}

\begin{interviewq}[$\star\star$ \iqroles{mle, researcher}]\label{iq:qm:linear-models-under-stress:4}
How do you cross-validate a model that predicts twenty-day returns from daily data?
\end{interviewq}

\begin{interviewq}[$\star\star$ \iqroles{researcher}]\label{iq:qm:linear-models-under-stress:5}
When are Fama--MacBeth \omterm{def:qm:estimation:estimator}{standard errors} wrong, and what do you use instead?
\end{interviewq}

\begin{interviewq}[$\star\star\star$ \iqroles{researcher, mle}]\label{iq:qm:linear-models-under-stress:6}
State and prove the Frisch--Waugh--Lovell theorem, and give a use of it on a trading desk.
\end{interviewq}
